\documentclass[addpoints]{exam}
\usepackage{amsmath}
\usepackage{amssymb}
\usepackage{color}
\usepackage{siunitx}

% \printanswers

\newcommand{\event}{Financial Mathematics}
\newcommand{\tournament}{}
\def\type{0}

\setlength\fillinlinelength{\textwidth}
\CorrectChoiceEmphasis{\color{red}\bfseries}
\SolutionEmphasis{\color{red}\bfseries}

\setlength\answerclearance{2pt}
\pagestyle{headandfoot}
\runningheadrule
\runningheader{\event}{\tournament}{\ifnum\type=1 Class Set Test \else Full Name:\hspace{1cm} \fi}
\runningfootrule
\runningfooter{}{Page \thepage\ of \numpages}{}

\begin{document}
\begin{center}
    {\huge Grade 12 Foundations for College Math \\ \event
    \ifprintanswers \space Key
    \else \space \ifnum\type<2 Test \fi \ifnum\type=1 \textbf{(CLASS SET)} \fi
          \ifnum\type=2 Answer Sheet \fi
    \fi \\ \vspace{3mm}\tournament\vspace{1mm}}
\end{center}

\vspace{.25\textheight}

\begin{center}
    First Name: \ifprintanswers
    \underline{\hspace{3.5cm}}\textbf{KEY}\underline{\hspace{5.5cm}}\\\ \vspace{5mm}
    \else \underline{\hspace{10cm}}\\ \vspace{5mm} \fi
    Last Name: \underline{\hspace{10cm}}\\ \vspace{5mm}
\end{center}

\noindent\textbf{\underline{Directions:}}
\begin{itemize}
    \ifnum\type=1 \item \textbf{This test is a class set.} Do not write on this paper. \fi
    \item Show all work for full marks. Round dollar amounts to the nearest cent.
    \item The test is designed to be completed in 75 minutes.
\end{itemize}
\begin{center}
\textit{For grading use only}\\
\multirowgradetable{1}[pages]
\end{center}

\newpage
\begin{questions}

\section*{Part A --- Multiple Choice (15 marks)}
\vspace{5mm}

% Q1
\question[1]
A principal of \$5{,}000 is invested at 4\% per year compounded annually for 3 years. The future value is closest to:
\begin{choices}
    \choice \$5{,}600.00
    \correctchoice \$5{,}624.32
    \choice \$5{,}612.00
    \choice \$5{,}648.00
\end{choices}
\vspace{3mm}

% Q2
\question[1]
Which formula correctly gives the future value $A$ of a principal $P$ compounded $n$ times per year at annual rate $r$ for $t$ years?
\begin{choices}
    \choice $A = P(1 + rt)$
    \correctchoice $A = P\!\left(1 + \dfrac{r}{n}\right)^{nt}$
    \choice $A = P\!\left(1 + \dfrac{r}{t}\right)^{n}$
    \choice $A = P e^{rt}$
\end{choices}
\vspace{3mm}

% Q3
\question[1]
An annuity pays \$500 at the end of each month for 10 years at 6\%/year compounded monthly. The interest rate per period $i$ is:
\begin{choices}
    \choice 6\%
    \choice 0.6\%
    \correctchoice 0.5\%
    \choice 1\%
\end{choices}
\vspace{3mm}

% Q4
\question[1]
The future value of an ordinary annuity formula is:
\begin{choices}
    \correctchoice $\displaystyle FV = PMT \times \frac{(1+i)^n - 1}{i}$
    \choice $\displaystyle FV = PMT \times \frac{1-(1+i)^{-n}}{i}$
    \choice $\displaystyle FV = PMT \times \frac{i}{(1+i)^n - 1}$
    \choice $\displaystyle FV = PMT \times (1+i)^n$
\end{choices}
\vspace{3mm}

% Q5
\question[1]
The present value of an ordinary annuity formula is:
\begin{choices}
    \choice $\displaystyle PV = PMT \times \frac{(1+i)^n - 1}{i}$
    \correctchoice $\displaystyle PV = PMT \times \frac{1-(1+i)^{-n}}{i}$
    \choice $\displaystyle PV = PMT \times i \cdot (1+i)^n$
    \choice $\displaystyle PV = \frac{PMT}{(1+i)^n}$
\end{choices}
\vspace{3mm}

% Q6
\question[1]
A lump sum of \$10{,}000 is invested at 5\% compounded semi-annually. The effective annual interest rate (EAR) is:
\begin{choices}
    \choice 5.00\%
    \correctchoice 5.0625\%
    \choice 5.125\%
    \choice 4.975\%
\end{choices}
\vspace{3mm}

% Q7
\question[1]
Which of the following describes an \textit{ordinary annuity}?
\begin{choices}
    \choice Payments made at the beginning of each period
    \correctchoice Payments made at the end of each period
    \choice A single lump-sum payment
    \choice Payments that grow by a fixed percentage each period
\end{choices}
\vspace{3mm}

% Q8
\question[1]
A mortgage has a principal of \$200{,}000 at 4\%/year compounded semi-annually. When calculating monthly payments, the equivalent monthly interest rate $i$ satisfies:
\begin{choices}
    \choice $i = 0.04/12$
    \correctchoice $i = (1.02)^{1/6} - 1$
    \choice $i = (1.04)^{1/12} - 1$
    \choice $i = 0.04/2$
\end{choices}
\vspace{3mm}

% Q9
\question[1]
A car loan of \$15{,}000 is taken at 6\%/year compounded monthly for 4 years. The number of payment periods $n$ is:
\begin{choices}
    \choice 4
    \choice 6
    \choice 36
    \correctchoice 48
\end{choices}
\vspace{3mm}

% Q10
\question[1]
Which financial goal is \textit{best} modelled by the present value of an annuity formula?
\begin{choices}
    \choice Saving a fixed amount each month to reach a target balance
    \correctchoice Determining how large a fund must be to support regular withdrawals
    \choice Calculating how much a savings account grows over time
    \choice Finding the compound interest on a single deposit
\end{choices}
\vspace{3mm}

% Q11
\question[1]
Total interest paid on a loan equals:
\begin{choices}
    \correctchoice Total of all payments minus the original principal
    \choice Original principal minus total of all payments
    \choice Monthly payment times annual rate
    \choice Principal times the number of years
\end{choices}
\vspace{3mm}

% Q12
\question[1]
An RRSP contribution of \$200/month is made for 30 years at 5\% compounded monthly. The total amount \textit{contributed} (not including interest) is:
\begin{choices}
    \choice \$60{,}000
    \correctchoice \$72{,}000
    \choice \$80{,}000
    \choice \$84{,}000
\end{choices}
\vspace{3mm}

% Q13
\question[1]
If you double the amortization period on a mortgage while keeping the interest rate and principal fixed, the monthly payment will:
\begin{choices}
    \choice Exactly halve
    \correctchoice Decrease, but not by exactly half
    \choice Stay the same
    \choice Increase
\end{choices}
\vspace{3mm}

% Q14
\question[1]
A \$1{,}000 bond pays 3\% simple interest per year for 5 years. The total amount received at maturity is:
\begin{choices}
    \choice \$1{,}159.27
    \choice \$1{,}030.00
    \correctchoice \$1{,}150.00
    \choice \$1{,}159.00
\end{choices}
\vspace{3mm}

% Q15
\question[1]
Two options for a \$25{,}000 loan: Option X charges 4\% compounded monthly over 5 years; Option Y charges 4\% compounded annually over 5 years. Which statement is true?
\begin{choices}
    \choice Option X and Option Y cost the same in total
    \correctchoice Option X costs more in total interest than Option Y
    \choice Option Y costs more in total interest than Option X
    \choice It depends on the monthly payment amount
\end{choices}
\vspace{3mm}

\section*{Part B --- Short Answer (33 marks)}

\vspace{5mm}

%% Q16 --- FV annuity / RRSP
\question[6]
\textbf{RRSP Savings Plan.}
Maya contributes \$300 at the end of each month into an RRSP that earns 6\%/year compounded monthly for 25 years.

\begin{parts}
\part[2] State the values of $PMT$, $i$, and $n$.
\part[2] Calculate the future value of Maya's RRSP at the end of 25 years.
\part[2] How much of that total is interest earned (not contributions)?
\end{parts}

\begin{solution}
\textbf{(a)} $PMT = \$300$; \quad $i = \dfrac{0.06}{12} = 0.005$; \quad $n = 25 \times 12 = 300$.

\textbf{(b)}
\[
FV = PMT \times \frac{(1+i)^n - 1}{i}
   = 300 \times \frac{(1.005)^{300} - 1}{0.005}
\]
$(1.005)^{300} = e^{300\ln(1.005)} \approx e^{300 \times 0.0049875} \approx e^{1.49625} \approx 4.4650$
\[
FV = 300 \times \frac{4.4650 - 1}{0.005}
   = 300 \times \frac{3.4650}{0.005}
   = 300 \times 693.00
   \approx \$207{,}900.94
\]
(Using a more precise value: $(1.005)^{300} \approx 4.46540$, giving $FV \approx \$207{,}924.07$.)

\textbf{(c)}
Total contributions $= 300 \times 300 = \$90{,}000$.\\
Interest earned $= \$207{,}924.07 - \$90{,}000 = \$117{,}924.07$.
\end{solution}

\vspace{5mm}

%% Q17 --- PV annuity / retirement fund
\question[6]
\textbf{Retirement Fund.}
Carlos wants to withdraw \$2{,}000 at the end of each month for 20 years from a retirement fund that earns 5\%/year compounded monthly.

\begin{parts}
\part[2] State the values of $PMT$, $i$, and $n$.
\part[2] Calculate the present value of the fund required.
\part[2] How much more than the total withdrawals is the required present value? (i.e., what is the interest portion working in Carlos's favour?)
\end{parts}

\begin{solution}
\textbf{(a)} $PMT = \$2{,}000$; \quad $i = \dfrac{0.05}{12} \approx 0.004\overline{16}$; \quad $n = 20 \times 12 = 240$.

\textbf{(b)}
\[
PV = PMT \times \frac{1-(1+i)^{-n}}{i}
   = 2000 \times \frac{1-(1.004\overline{16})^{-240}}{0.004\overline{16}}
\]
$(1.004\overline{16})^{240} \approx 2.7126$, so $(1.004\overline{16})^{-240} \approx 0.36861$.
\[
PV = 2000 \times \frac{1 - 0.36861}{0.004\overline{16}}
   = 2000 \times \frac{0.63139}{0.004\overline{16}}
   \approx 2000 \times 151.525
   \approx \$303{,}049.
\]

\textbf{(c)}
Total withdrawals $= 2000 \times 240 = \$480{,}000$.\\
Difference $= \$480{,}000 - \$303{,}049 = \$176{,}951$.\\
The fund earns approximately \$176{,}951 in interest over 20 years, which is what allows Carlos to withdraw more than he started with.
\end{solution}

\vspace{5mm}

%% Q18 --- Mortgage
\question[5]
\textbf{Mortgage Payment.}
A homebuyer takes a \$350{,}000 mortgage at 4.8\%/year compounded semi-annually, amortized over 25 years with monthly payments.

\begin{parts}
\part[2] Convert the semi-annual rate to an equivalent monthly rate $i$.
\part[1] State the number of payment periods $n$.
\part[2] Find the monthly payment and the total interest paid over 25 years.
\end{parts}

\begin{solution}
\textbf{(a)} The nominal rate is 4.8\%/year compounded semi-annually, so the semi-annual rate is $\dfrac{0.048}{2} = 0.024$.

Equivalent monthly rate:
\[
(1+i)^{12} = (1.024)^{2} \implies i = (1.024)^{2/12} - 1 = (1.024)^{1/6} - 1
\]
\[
i = (1.024)^{0.1\overline{6}} - 1 \approx 1.003965 - 1 = 0.003965 \quad (0.3965\%/\text{month})
\]

\textbf{(b)} $n = 25 \times 12 = 300$ monthly payments.

\textbf{(c)} Using the PV annuity formula rearranged for PMT:
\[
PMT = PV \times \frac{i}{1-(1+i)^{-n}}
    = 350{,}000 \times \frac{0.003965}{1-(1.003965)^{-300}}
\]
$(1.003965)^{300} \approx e^{300 \times 0.003957} \approx e^{1.1871} \approx 3.2779$, so $(1.003965)^{-300} \approx 0.30507$.
\[
PMT = 350{,}000 \times \frac{0.003965}{1 - 0.30507}
    = 350{,}000 \times \frac{0.003965}{0.69493}
    \approx 350{,}000 \times 0.005705
    \approx \$1{,}996.63
\]
Total paid $= 1{,}996.63 \times 300 \approx \$598{,}989$.\\
Total interest $= \$598{,}989 - \$350{,}000 = \$248{,}989$.
\end{solution}

\vspace{5mm}

%% Q19 --- Loan comparison
\question[8]
\textbf{Car Loan Comparison.}
Priya needs to borrow \$20{,}000 to buy a car and is comparing two financing options:
\begin{itemize}
    \item \textbf{Option A:} 5\%/year compounded monthly, repaid over 5 years
    \item \textbf{Option B:} 7\%/year compounded monthly, repaid over 4 years
\end{itemize}

\begin{parts}
\part[2] Find the monthly payment for Option A.
\part[2] Find the monthly payment for Option B.
\part[2] Find the total amount paid under each option.
\part[2] Which option do you recommend, and why? Consider both the monthly payment and the total cost.
\end{parts}

\begin{solution}
\textbf{(a) Option A:} $PV = 20{,}000$, $i = 0.05/12 \approx 0.004167$, $n = 60$.
\[
PMT_A = 20{,}000 \times \frac{0.004167}{1-(1.004167)^{-60}}
\]
$(1.004167)^{60} \approx 1.28336$, so $(1.004167)^{-60} \approx 0.77921$.
\[
PMT_A = 20{,}000 \times \frac{0.004167}{0.22079} \approx 20{,}000 \times 0.018871 \approx \$377.42
\]

\textbf{(b) Option B:} $PV = 20{,}000$, $i = 0.07/12 \approx 0.005833$, $n = 48$.
\[
PMT_B = 20{,}000 \times \frac{0.005833}{1-(1.005833)^{-48}}
\]
$(1.005833)^{48} \approx 1.32290$, so $(1.005833)^{-48} \approx 0.75590$.
\[
PMT_B = 20{,}000 \times \frac{0.005833}{0.24410} \approx 20{,}000 \times 0.023893 \approx \$477.86
\]

\textbf{(c)}
Total A $= 377.42 \times 60 = \$22{,}645.20$; \quad interest A $= \$2{,}645.20$.\\
Total B $= 477.86 \times 48 = \$22{,}937.28$; \quad interest B $= \$2{,}937.28$.

\textbf{(d)} Option A is recommended if the higher total cost of Option B matters more. Option A has a lower monthly payment (\$377 vs \$478) and a lower total cost (\$22{,}645 vs \$22{,}937). Option B finishes one year sooner but costs more overall. For most budgets, Option A is preferable.
\end{solution}

\vspace{5mm}

%% Q20 --- Retirement planning (FV + withdrawal)
\question[8]
\textbf{Retirement Planning.}
At age 30, Dani begins investing \$500 at the end of each month into a retirement account that earns 7\%/year compounded monthly. She plans to retire at age 65.

\begin{parts}
\part[2] How many payments does Dani make?
\part[3] Calculate the future value of her account at retirement.
\part[3] After retirement, Dani wants equal monthly withdrawals over 25 years from the same account (still earning 7\%/year compounded monthly). Find the monthly withdrawal amount.
\end{parts}

\begin{solution}
\textbf{(a)} Number of payments: $n = (65 - 30) \times 12 = 35 \times 12 = 420$.

\textbf{(b)}
$PMT = 500$, $i = 0.07/12 \approx 0.005833$, $n = 420$.
\[
FV = 500 \times \frac{(1.005833)^{420} - 1}{0.005833}
\]
$(1.005833)^{420} = e^{420 \times \ln(1.005833)} \approx e^{420 \times 0.005817} \approx e^{2.4431} \approx 11.509$.
\[
FV = 500 \times \frac{11.509 - 1}{0.005833}
   = 500 \times \frac{10.509}{0.005833}
   \approx 500 \times 1{,}801.7
   \approx \$900{,}854
\]

\textbf{(c)}
Now $PV = \$900{,}854$ (the retirement balance), $i = 0.005833$, $n = 25 \times 12 = 300$.
\[
PMT = 900{,}854 \times \frac{0.005833}{1-(1.005833)^{-300}}
\]
$(1.005833)^{300} \approx e^{300 \times 0.005817} \approx e^{1.7451} \approx 5.727$, so $(1.005833)^{-300} \approx 0.17461$.
\[
PMT = 900{,}854 \times \frac{0.005833}{0.82539}
    \approx 900{,}854 \times 0.007067
    \approx \$6{,}367\ \text{per month}
\]
\end{solution}

\end{questions}
\end{document}
